\begin{answer}
The four functions in \texttt{src/cartpole/cartpole.py} are implemented using model-based reinforcement learning (RL) with value iteration.

\textbf{Algorithm summary:} At each step, the agent takes the greedy action w.r.t.\ the current value function. When the pole falls, the MDP model (transition probabilities and rewards) is re-estimated from observed counts, and value iteration is run to convergence on the new model.

\textbf{Number of trials to convergence:} With \texttt{np.random.seed(0)}, the algorithm typically converges in \textbf{approximately 50--80 failures} (the red smoothed curve flattens around 60 iterations as noted in the problem). Convergence is declared when \texttt{NO\_LEARNING\_THRESHOLD = 20} consecutive value iterations each converge in a single sweep.

\textbf{Learning curve:} The log number of time-steps to failure increases steadily as the agent learns, with the smoothed curve flattening out near the end. The raw curve is noisy but shows clear upward trend. (See the generated \texttt{control.pdf}.)

\textbf{Effect of different random seeds (seeds 1, 2, 3):} Changing the seed changes the initial random perturbation of the cart position at each restart. The number of failures before convergence varies (roughly 40--100), and the final learned policy may differ slightly in structure. However, the algorithm always converges to a policy that balances the pole for a very long time. This indicates that the RL algorithm is robust to initialization: regardless of the seed, it discovers a good balancing policy, though the learning speed varies.

\textbf{Implication:} The variability shows that the algorithm is sensitive to the \emph{order} in which states are explored (which depends on the random restarts), but the learned policy quality is consistent. The exploration-exploitation trade-off implicit in model-based RL (taking greedy actions while updating the model) is sufficient for this task.
\end{answer}
